\section{A strict segment inequality}\label{sec:08-segments}

We use the constants and profiles of Appendix~\ref{sec:07-scalar}.
In particular, all finite decimals below are exact rationals. For an interval
$[x_-,x_+]$, a bar denotes its uniform average in the physical coordinate $x$;
endpoint subscripts are evaluations, not derivatives. The coordinate
$t=\operatorname{arsinh}(x/s)$ and the notation
$n_F=\ddot F-\tanh(t)\dot F=(s\cosh t)^2F''(x)$ are as in that Section.
When a profile is written below with a $t$-coordinate argument, it
means its composition with $x=s\sinh t$; for example $\mathsf K(t)$
means $\mathsf K(s\sinh t)$. A subscript $+$ or $-$ still denotes
its value at the corresponding physical endpoint $x_+$ or $x_-$.
Proposition~\ref{prop:segment-estimate} proves the segment estimate used in the matrix
argument.

\begin{proposition}[Strict segment inequality]\label{prop:segment-estimate}
Put $u=(\pi,q)$. For every pair $x_-<x_+$, define $e_K$ and $\Gamma$ as below.
Then
\[
\begin{aligned}
 e_K&+\Gamma+
 \frac{\sum_{\sigma\in\{-,+\}}(\overline{\mathsf C}-\mathsf C_{-\sigma}
                     -\overline{|u-u_\sigma|^2})^2}{8\lambda}
 +t_+\,\frac12\sum_{\sigma\in\{-,+\}}|\overline u-u_\sigma|^2
 < \overline{|u|^2}-|\overline u|^2,\\
 e_K&=\overline{\mathsf K}-(\mathsf K_++\mathsf K_-)/2,\\
 \Gamma&=\frac{.80}{m_*^2}
 \left(1.25(\mathsf K_+-\mathsf K_-+t_c(\mathsf C_+-\mathsf C_-))^2+
             5t_r^2(\mathsf C_+-\mathsf C_-)^2\right),
\end{aligned}
\]
where $-\sigma$ denotes the opposite endpoint. At coincident endpoints,
both sides of the asserted inequality are zero.
\end{proposition}

\subsection{Coordinates and reduction to scalar budgets}
\label{subsec:segments-reduction}

We prove Proposition~\ref{prop:segment-estimate} by separating intervals according to their width
and location in the $t$ coordinate. Write the endpoint coordinates as
\[
 l=m-h,\qquad y=m+h,\qquad h>0.
\]
Here $m$ is a scalar midpoint, unrelated to the ambient dimension.
The profiles $\mathsf K,\mathsf C,q$ are even, while reflection applies
the fixed orthogonal map $(u_1,u_2,u_3)\mapsto(u_1,-u_2,u_3)$ to $u$.
It exchanges the endpoints and preserves every average and squared
quantity in Proposition~\ref{prop:segment-estimate}. We may therefore assume $m\ge0$, so that
$y\ge|l|$.

The estimates below cover the endpoint pairs in the following order.
This division leaves the final coverage check, including all shared
boundaries, to the end of the proof.
\[
\begin{array}{ll}
 0<h\le .16 & \text{short widths, Subsection~\ref{subsec:segments-short}},\\
 h\ge .16\text{ and a tail case} &
       \text{Subsection~\ref{subsec:segments-tails}},\\
 |l|\le1.4,\ y\le3.6,\ .16\le h\le.54 &
       \text{middle widths, Subsection~\ref{subsec:segments-middle}},\\
 |l|\le1.4,\ y\le3.6,\ h\ge.54 &
       \text{larger widths, Subsection~\ref{subsec:segments-large}}.
\end{array}
\]
The tail cases mean $l\le-1.4$, $l\ge1.4$, or $|l|<1.4$ with
$y\ge3.6$. The remaining bounded region uses the auxiliary estimates
in Subsection~\ref{subsec:segments-auxiliary}.

Since $dx=s\cosh(t)\,dt$, the normalized averaging density on $[l,y]$ is
\[
 \frac{\cosh t}{2\cosh m\sinh h}.
\]
Identifying the two circle coordinates with a complex number and integrating
$re^{iwt}\cosh t$ gives the mean
$re^{iwm}(A_h+iB_h\tanh m)$, where
\[
 A_h=\frac{\cos(wh)+w\coth h\sin(wh)}{1+w^2},\qquad
 B_h=\frac{\coth h\sin(wh)-w\cos(wh)}{1+w^2}.
\]
Put
\[
\begin{aligned}
 f_0^*&=\frac{r^2}{1+w^2}\left(w^2-\frac{\sin^2(wh)}{\sinh^2 h}\right),
 &p_a&=2r^2(1-A_h\cos(wh)),& p_d&=-2r^2 B_h\tanh m\sin(wh),\\
 b_\sigma&=\overline{(q-q_\sigma)^2},
 &b_a&=(b_++b_-)/2,& b_d&=(b_+-b_-)/2,\\
 &&c_a&=\overline{\mathsf C}-(\mathsf C_++\mathsf C_-)/2,& c_d&=(\mathsf C_+-\mathsf C_-)/2 .
\end{aligned}
\]
The elementary identity
\[
 A_h^2+B_h^2
 =\frac{1+\sin^2(wh)/\sinh^2h}{1+w^2}
\]
shows that $f_0^*=r^2(1-A_h^2-B_h^2)$ lower bounds the variance of the
circle coordinates. It is strictly positive for $h>0$, because
$|\sin(wh)|\le wh<w\sinh h$. Taking scalar products with the two
endpoint vectors also gives
$\overline{|\pi-\pi_\pm|^2}=p_a\pm p_d$. Consequently
\[
 \overline{\mathsf C}-\mathsf C_{\mp}-\overline{|u-u_\pm|^2}
             =(c_a-p_a-b_a)\pm(c_d-p_d-b_d),\qquad
 \tfrac12\sum_\sigma|\bar u-u_\sigma|^2
             =p_a+b_a-(\overline{|u|^2}-|\bar u|^2).
\]
Let $V_u=\overline{|u|^2}-|\bar u|^2$ and
$V_q=\overline{q^2}-\bar q^2$. The two preceding identities turn the
inequality of Proposition~\ref{prop:segment-estimate} into the assertion that its squared terms plus
$e_K+\Gamma+t_+(p_a+b_a)$ are less than $(1+t_+)V_u$.
Because $V_u\ge f_0^*+V_q$ and $1+t_+>0$, it will usually suffice to prove
\[
 e_K+\Gamma+
 \frac{(c_a-p_a-b_a)^2+(c_d-p_d-b_d)^2}{4\lambda}
 +t_+(p_a+b_a)<(1+t_+) f_0^*. \tag{s0}\label{eq:s0}
\]

\begin{lemma}[Chord kernel]\label{lem:segment-chord}
For every twice continuously differentiable profile $F$ on the segment,
let $t=m+h\beta$ and
\[
 \theta=\frac{\sinh t-\sinh(m-h)}{2\cosh m\sinh h},\qquad
 S=\frac{\sinh h}{h},\qquad g_h=\frac34S(1+\cosh h).
\]
Then
\[
 \overline F-(F_-+F_+)/2 =-\frac{h^2}{3}\int_{-1}^1 W(\beta)\,n_F(m+h\beta)\,\frac{d\beta}2,\qquad
 W=6\theta(1-\theta)\frac{\cosh m}{\cosh t}\frac{\sinh h}{h},
\]
where $n_F$ is evaluated at $x=s\sinh(m+h\beta)$. Moreover,
\[
 0\le W\le\frac32S,\qquad
 W\le g_h(1-\beta^2),\qquad
 \int_{-1}^1W\,\frac{d\beta}{2}\le\frac S2(1+\cosh h).
\]
\end{lemma}
\begin{proof}
The uniform trapezoidal error in the $x$ coordinate is
$-(x_+-x_-)^2\int_0^1\theta(1-\theta)F''(x_-+\theta(x_+-x_-))\,d\theta/2$.
Substitute $x_+-x_-=2s\cosh m\sinh h$ and
$d\theta=h\cosh t\,d\beta/(2\cosh m\sinh h)$ to obtain the stated identity.
For the first upper bound on $W$, rearranging its definition reduces the
claim to
$(\sinh t-\sinh m\cosh h)^2\ge
\cosh m(\cosh m-\cosh t)\sinh^2h$.
To check this inequality, regard $m,t$ as fixed and put
$\zeta=\cosh h-1\ge0$. The difference of its two sides is
\[
 (\sinh t-\sinh m)^2
 +2(\cosh(m-t)-1)\zeta
 +(\cosh t\cosh m-1)\zeta^2.
\]
Every coefficient is nonnegative. This proves $W\le3S/2$.

The sharper bound near the endpoints comes from factoring the two
differences of hyperbolic sines in $\theta(1-\theta)$. With
$u_\pm=(1\pm\beta)/2$, the factorization gives
\[
 \theta(1-\theta)\frac{\cosh m}{\cosh t}
 =\frac{\sinh(u_+h)\sinh(u_-h)}{\sinh^2h}
  \frac{\cosh(2m+\beta h)+\cosh h}
       {\cosh(2m+\beta h)+\cosh(\beta h)}.
\]
Since $0\le u_\pm\le1$, convexity gives
$\sinh(u_\pm h)\le u_\pm\sinh h$. Thus the product on the left is bounded by
\[
 \frac{1-\beta^2}{4}\frac{\cosh(2m+\beta h)+\cosh h}{\cosh(2m+\beta h)+\cosh(\beta h)}
\]
The quotient in this display is at most $(1+\cosh h)/2$: subtract one
and use $\cosh(2m+\beta h)+\cosh(\beta h)\ge2$.
This proves the pointwise bound involving $g_h$. Its integral uses
$\int_{-1}^1(1-\beta^2)\,d\beta/2=2/3$ and gives
$\int W\,d\beta/2\le S(1+(\cosh h-1)/2)$, as required.
\end{proof}

\subsection{Short widths}\label{subsec:segments-short}

Suppose $h\le.16$. Lemma~\ref{lem:segment-chord} bounds $W$ by
$1.507$ and its integral by $1.012$. Put $G=r^2w^2$ and
$u_0=Gh^2/3$, the scale of the circle variance on a short interval.
By the chord identity and Cauchy--Schwarz applied to the
endpoint derivative integrals, the joint bound $U_*$ in
Lemma~\ref{lem:scalar-ranges} gives
\[
 e_K+\Gamma
 \le\frac{h^2}{3}\int_{-1}^1
 U_*(m+h\beta)\,\frac{d\beta}{2}
 \le .54u_0.
\]
The coefficient $12\cdot .80/m_*^2$ in $U_*$ arises from
$(F_+-F_-)^2\le4h^2\int_{-1}^1\dot F(m+h\beta)^2\,d\beta/2$.
The other terms obey
\[
\begin{array}{ll}
 f_0^*/u_0\ge (1-2w^4 h^2/(15(1+w^2)))/(1+h^2/3),
 & |p_d|\le 2h u_0,\\
 p_a/u_0\le4+.15 h^2,
 & b_a/u_0\le(.193)^2(4+.15h^2)/G,\\
 |c_a|/u_0\le.79\cdot1.012/G,& |c_d|\le.247 h,\quad
 |b_d|\le2(.193)h\cdot .53\cdot1.012h^2/3 .
\end{array}
\]
We justify the table by separating the angular estimates from the
endpoint and chord errors. For the lower bound on $f_0^*$, substitute
$\sinh^2h\ge h^2+h^4/3$ and
$\sin^2 y\le y^2-y^4/3+2y^6/45$, valid for $|y|\le2$, in its
defining formula. For the angular half-difference, write the numerator
of $B_h$ as
$(\sin(wh)-wh\cos(wh))/h+(\coth h-1/h)\sin(wh)$.
The estimates $|\sin y-y\cos y|\le|y|^3/3$ and
$0\le\coth h-1/h\le h/3$ give $|B_h|\le wh^2/3$, hence
$|p_d|\le2hu_0$. The latter hyperbolic bound follows by comparing
the coefficients of $h\cosh h$ and $(1+h^2/3)\sinh h$.

For the endpoint distances, the derivatives are taken in $t$, but the
expectation is still uniform in physical length. Odd terms cancel in
the second moment of $\beta=(t-m)/h$, so
\[
 \overline{\beta^2}
 =\frac{\int_{-1}^1\beta^2\cosh(h\beta)\,d\beta/2}{S}
 \le\frac{1/3+h^2/10+h^4\cosh h/168}{1+h^2/6}
 \le\frac13+.05h^2\qquad(h\le.54).
\]
The first inequality uses the first three even terms and their Taylor
remainder. For the last, $\cosh h\le(1-h^2/2)^{-1}<7/5$ on this
range, and cross-multiplication gives the stated bound.
If a profile has $t$-derivative norm at most $L$, then the half-sum of
its squared distances to the two endpoint values is bounded, after
averaging, by $L^2h^2(1+\overline{\beta^2})$.
Apply this with $L=rw$ for the circle and $L=.193$ for $q$.
Division by $u_0$ gives the bounds on $p_a$ and $b_a$, including their
factor $4$.

For the $\mathsf C$ terms, the chord identity gives
$|c_a|\le.79\cdot1.012h^2/3$, and integration of
$|\dot{\mathsf C}|\le.247$ gives $|c_d|\le.247h$.
Finally,
$|b_d|=|q_+-q_-|\,|\bar q-(q_++q_-)/2|$.
The derivative bound controls the first factor by $2(.193)h$;
the chord identity controls the second by $.53\cdot1.012h^2/3$.
These give the last entries of the table.

It remains to compare these costs with the right side of
Equation~\eqref{eq:s0}. Subtract the two squared terms and the
$t_+$ term from that right side, then divide by $u_0$.
The preceding bounds give the following lower bound, with
$H=.16$, $P=4+.15H^2$, and $Q=.193^2P/G$:
\[
\begin{aligned}
 &(1+t_+)\frac{1-2w^4H^2/(15(1+w^2))}{1+H^2/3}-t_+(P+Q)\\
 &\quad-\frac1{4\lambda}\Bigl[\frac{GH^2}{3}(P+Q+.79\cdot1.012/G)^2\\
 &\hspace{95pt}+\frac3G\bigl(.247+\tfrac23(G+.193\cdot.53\cdot1.012)H^2\bigr)^2\Bigr].
\end{aligned}
\]

The displayed rational lower bound is greater than $.55510$, hence greater
than $.55$. It exceeds the normalized cost $.54$, proving
Equation~\eqref{eq:s0} throughout this short-width range.

\subsection{Segments with endpoints in the tails}\label{subsec:segments-tails}

We next handle all segments with $h\ge.16$ that leave the central region.
All profile ranges and positive-part integral bounds in this subsection
come from Lemmas~\ref{lem:scalar-ranges} and~\ref{lem:scalar-tail}.

\paragraph{A central endpoint and a distant endpoint.}
If \(|m-h|<1.4,\ m+h\ge3.6\), then \(h\ge1.1\) and we can use
\[
 f_0^*\ge.0449,\quad .0818\le p_a\le.1074,\quad |p_d|\le.0132 .
\]
\paragraph{Endpoints in opposite tails.}
If $m-h\le-1.4$, then \(h\ge1.4,\ m+h\ge1.4\), and bounds here are
\[
 f_0^*\ge.0456,\quad .0830\le p_a\le.1062,\quad |p_d|\le.0129.
\]
These angular bounds follow by diagonalizing the two elementary
quadratic forms in $\cos(wh)$ and $\sin(wh)$:
\[
\begin{aligned}
 \frac{1-\sqrt{1+w^2\coth^2h}}{2(1+w^2)}
 &\le A_h\cos(wh)\le
 \frac{1+\sqrt{1+w^2\coth^2h}}{2(1+w^2)},\\
 |B_h\sin(wh)|&\le
 \frac{\coth h+\sqrt{\coth^2h+w^2}}{2(1+w^2)}.
\end{aligned}
\]
Here $\sinh h>1.335,1.904$ and $\coth h<1.25,1.13$ in the two
cases, respectively. The positive-term series for $\sinh h$ and
$\exp(2h)$ give these bounds.

In the central/distant case, the physical endpoints satisfy
$x_+>65.8$ and $|x_-|<6.86$. If the segment contains zero, its length
is at least $65.8$. The integral of $(\mathsf K)_+$ over its positive
part is at most $.422$, and over its negative part at most $.126$.
If the segment does not contain zero, its length is at least
$65.8-6.86$ and the integral is at most $.422$. Therefore
\[
 \overline{\mathsf K}
 \le\max\left\{\frac{.422+.126}{65.8},
                    \frac{.422}{65.8-6.86}\right\}<.0084.
\]
For $b_+$, separate the physical set $|x|<5$ from its complement.
Its probability is at most $10/(65.8-6.86)$ under the uniform
segment measure. The far endpoint has $q_+\in[.255,.2556]$ by
Lemma~\ref{lem:scalar-tail}, since $x_+>64$. Globally
$q\in[.0775,.2559]$, and on the complement $q\in[.2203,.2559]$.
Consequently
\[
 b_+\le\frac{10}{65.8-6.86}(.1781)^2
       +\left(1-\frac{10}{65.8-6.86}\right)(.0353)^2<.0068.
\]
The global oscillation also gives $b_-\le.1784^2<.032$.
The one-variable $\mathsf C$ ranges give
$c_a\in[-.081,.160]$ and $|c_d|\le.080$.

In the opposite-tail case, both endpoints have $|t|\ge1.4$.
Use the global bound $\overline{\mathsf K}\le.0275$ and
$b_\pm\le.1784^2<.032$. The endpoint $\mathsf C$ range is now
$[-.501,-.495]$, giving $-.006\le c_a\le.160$ and
$|c_d|\le.003$.

For clarity, to bound \(e_K+\Gamma\) in these cases take intervals \(\mathcal I_-,\mathcal I_+\) for \(\mathsf K_-,\mathsf K_+\) and \(o=|\mathsf C_+-\mathsf C_-|_{\max}\), from the one-variable bounds:
\[
 \begin{array}{c|ccc|cc}
 &\mathcal I_- &\mathcal I_+&o&\overline{\mathsf K}_{\max}
 &\text{upper bound on }e_K+\Gamma\\\hline
 \text{central/distant} &[-.007,.0275]&[0,.0008]&.160&.0084&.0144\\
 \text{opposite tails} &[-.0001,.0191]&[-.0001,.0191]&.006&.0275&.0277
 \end{array}
\]
After replacing $|\mathsf C_+-\mathsf C_-|$ by $o$, the expression to
maximize is convex in the two $\mathsf K$ endpoint values. Its maximum
on the rectangle therefore occurs at a vertex. Add to
$\overline{\mathsf K}_{\max}$ the maximum, over pairs of endpoints
$z_\pm$ of these intervals, of
$-(z_++z_-)/2+.80[1.25(|z_+-z_-|+t_co)^2+5t_r^2o^2]/m_*^2$.
This gives the last column of the table.

For the other costs in Equation~\eqref{eq:s0}, $|b_d|\le.016$ in
both cases. The central/distant bounds give
$b_a\le.0194$, $|c_a-p_a-b_a|\le.2078$, and
$|c_d-p_d-b_d|\le.1092$. Thus the remaining budget is at least
\[
 1.064(.0449)-\frac{.2078^2+.1092^2}{2.6}
              -.064(.1074+.0194)>.018>.0144.
\]
For opposite tails, the corresponding bounds are $b_a\le.032$,
$|c_a-p_a-b_a|\le.1442$, and $|c_d-p_d-b_d|\le.0319$.
They give
\[
 1.064(.0456)-\frac{.1442^2+.0319^2}{2.6}
              -.064(.1062+.032)>.031>.0277.
\]
In each case the remaining budget exceeds the bound on $e_K+\Gamma$.

\paragraph{Both endpoints in the positive tail.}
If $m-h\ge1.4$ and $h\ge.16$, then \(e_K\le0\) by convexity, and
\[
 \Gamma/f_0^*\le .10,\qquad
 p_a+|p_d|\le \min(.150,\,6.65 f_0^*),\qquad f_0^*\ge.008 .
\]
To bound $\Gamma/f_0^*$, first suppose $h\le.4$.
The short-width lower bound on $f_0^*/u_0$ remains valid up to
$wh=2$, and gives $f_0^*/h^2>.18$ here. Combine it with the
derivative bounds on the endpoint differences in $\Gamma$.
For $h\ge.4$, instead use $f_0^*\ge.0332$, following from
$\sinh(.4)>.410$, and the value ranges of the endpoint profiles.
Both estimates give $\Gamma/f_0^*\le.10$.

For the circle endpoint distance, when $h\le.3$ the short-width
estimates give $p_a+|p_d|\le(4+.15h^2+2h)u_0$.
Together with the same variance lower bound this gives both asserted
upper bounds on $p_a+|p_d|$. For $h\ge.3$, use the angular
quadratic-form bounds with $\coth(.3)<3.434$ to obtain $.150$.
The variance bound $f_0^*\ge.02256$, from $\sinh(.3)>.304$,
then gives $.150<6.65f_0^*$.
On the remaining range $[.16,.3]$, the functions $\sin(wh)/h$
and $h/\sinh h$ are positive decreasing. Thus $f_0^*$ is increasing,
and its lower bound at $.16$ gives $f_0^*\ge.008$ throughout.

All points of this segment lie in the positive tail. The profile
ranges therefore give $|\overline{\mathsf C}-\mathsf C_\pm|\le.006$
and $b_\pm\le.00045$.
The sum of the two squares in Equation~\eqref{eq:s0} is the mean,
over endpoint labels $\sigma$, of
$(\overline{\mathsf C}-\mathsf C_{-\sigma}
-\overline{|u-u_\sigma|^2})^2$.
Each circle contribution is bounded both by $.150$ and by
$6.65f_0^*$, hence by $\sqrt{.150\cdot6.65f_0^*}$.
After division by $f_0^*$, the whole left side of
Equation~\eqref{eq:s0} is therefore at most
\[
 .10+\frac1{4\lambda}\left(\sqrt{.150\cdot6.65}+\frac{.00645}{\sqrt{.008}}\right)^2+
 t_+(6.65+.00045/.008)<1+t_+ .
\]

\subsection{Auxiliary estimates in the bounded region}
\label{subsec:segments-auxiliary}

The preceding cases leave \(h\ge .16\), \(l=m-h\in[-1.4,1.4]\), \(y=m+h\le3.6\), with $m\ge0$ and $y\ge|l|$.
For middle widths with $l\ge0$, retaining the $q$ variance will offset
part of the endpoint-distance cost $b_a$. We therefore use the sufficient
inequality
\[
 e_K+\Gamma+\big[(c_a-p_a-b_a)^2+(c_d-p_d-b_d)^2\big]/(4\lambda)+t_+(p_a+b_a)
     < (1+t_+)(f_0^*+V_q).                                                   \tag{sv}\label{eq:sv}
\]
Indeed the squared deviations of the mean at the endpoints average to \(p_a+b_a-(\overline{|u|^2}-|\bar u|^2)\).

Put
\[
\begin{gathered}
 \mathsf D=\mathsf K+t_c\mathsf C,\qquad
 d_*^K=\frac{.80\cdot1.25}{m_*^2},\qquad
 d_*^C=\frac{.80\cdot5t_r^2}{m_*^2},\\
 T_0=.54,\qquad S=\frac{\sinh h}{h},\qquad
 g_h=\frac34S(1+\cosh h).
\end{gathered}
\]
Thus \(\Gamma=d_*^K(\mathsf D(y)-\mathsf D(l))^2+d_*^C(\mathsf C(y)-\mathsf C(l))^2\).
Here are concentration constants:
\[
\begin{array}{c|c|cc}
 j & A_j(t)& \alpha_j&\beta_j\\\hline
 D&\dot{\mathsf D}&.07052&.0754\\
 E&-\dot{\mathsf C}&.281&.254\\
 N&-n_{\mathsf K}&.24&.277\\
 C& n_{\mathsf C}&.728&.812
\end{array}
\qquad L_j(H)=2\alpha_j/3-\beta_j H/2.
\]
\begin{lemma}[Concentration and primitive bounds]\label{lem:segment-auxiliary}
The following bounds hold in the bounded region just specified. For a set in \([0,3.6]\) of length \(2H\), \(0<H\le T_0\), the mean of \((A_j)_+\) on the set is at most \(\alpha_j-\beta_j H\). Also \(\dot{\mathsf D}\ge-.022\) on \([0,3.6]\), and for \(0\le t\le t'\le3.6\),
\[
 \mathsf C(t')-\mathsf C(t)\le .000124,\qquad q(t)-q(t')\le .000084 .
\]
We will use \(b_a\le .0163\), and primitive bounds
\[
\begin{array}{ll}
 l<0: & e_K+\Gamma\le \min(.018,.0146\cdot4 h^2+.00010)+.00015,\\
 l\ge0: & e_K+\Gamma\le .01815+1/(312\sinh^2(2h)).
\end{array}                                                               \tag{s1}\label{eq:s1}
\]

\end{lemma}
\begin{proof}
We establish the concentration, interpolation, and primitive assertions in turn.

\smallskip
\noindent\textbf{Concentration from finite threshold sums.}\quad
For the concentration assertion, compare with a threshold \(v\): the total on any given set is at most \(2Hv+I_j(v)\), where \(I_j(v)=\int_0^{3.6}(A_j-v)_+ dt\) and we use \(v\ge0\). The same threshold bound holds for a measurable
density $0\le\vartheta\le1$ on $[0,3.6]$ with $\int_0^{3.6}\vartheta=2H$:
\[
 \int_0^{3.6}\vartheta(t)(A_j(t))_+\,dt
 \le 2Hv+I_j(v),\qquad v\ge0.
\]
This follows by multiplying the pointwise threshold bound by $\vartheta$.

Here is a small table for this comparison on successive intervals of \(1000H\) with endpoints
\[
 0,75,150,200,250,300,350,400,540.
\]
Row entries give \(1000v\), followed below by upper bounds for \(10^5 I_j(v)\) in the same order.
\[
\begin{array}{c|rrrrrrrr}
 D&61&54&44&37&29&21&14&0\\
 E&247&224&192&167&141&116&90&42\\
 N&216&178&143&115&77&42&17&0\\
 C&664&545&444&347&234&152&97&21\\\hline
 D&0&105&419&716&1124&1598&2069&3167\\
 E&0&351&1326&2363&3673&5133&6885&10779\\
 N&0&577&1576&2624&4378&6352&8026&9323\\
 C&0&1776&4528&8066&13250&17831&21389&27670
\end{array}
\]
At both ends of each interval these upper bounds divided by \(10^5\) are \(\le 2H(\alpha_j-\beta_j H-v)\), which suffices by concavity. For the integral entries use exactly
\[
 \sum_{i=0}^{144}\frac{\delta_i}{80}(A_j(t_i)+\varepsilon_j-v)_+,\qquad
 \varepsilon_j=.00040,\ .00120,\ .00217,\ .00741
\]
respectively, with \(t_i=i/40\), \(\delta_i=2\) except \(\delta_0=\delta_{144}=1\), using the values and interval rules of the one-variable estimates. Indeed by the chord errors there, each \(A_j\) has upper error from its chord at most \(\varepsilon_j\), so the formula works by convexity. The sample values are enclosed by the finite rational rules of
Lemma~\ref{lem:scalar-certificates}; the interpolation estimates in
Lemma~\ref{lem:scalar-interpolation} justify the passage between nodes.
For example, with $\mathcal R(M)=.536M+.00001$, the required errors are
\[
\begin{aligned}
 \mathcal R(.000673)+.021\mathcal R(.00222)&=.00039592632<.00040,\\
 \mathcal R(.00222)&=.00119992<.00120,\\
 \mathcal R(.00402)&=.00216472<.00217,\\
 \mathcal R(.0138)&=.0074068<.00741.
\end{aligned}
\]
Thus these sums are rigorous upper bounds for the integrals; no
additional derivative evaluation between nodes is needed.

\smallskip
\noindent\textbf{Averaging the interpolated profiles.}\quad
For the remaining assertions put
\[
 P_F(t)=\int_0^1 F(ts')\cosh(ts')\,ds'\ \Big/\int_0^1\cosh(ts')\,ds',
 \qquad P=P_{\mathsf K}.
\]
Use as approximate samples (starred) at \(t_i\), with \(F_i=F(t_i)\),
\[
 P_{F,i}^*=F_i-40\sum_{k=1}^i\frac{\cosh t_k-\cosh t_{k-1}}{\sinh t_i}(F_k-F_{k-1}),
 \qquad P_{F,0}^*=F_0.
\]
This is an exact integration formula for the piecewise affine interpolant
of $F$ in the $t$ coordinate, averaged against physical length. Indeed,
integration by parts gives
$P_F(t)=F(t)-\int_0^t\dot F(v)\sinh v\,dv/\sinh t$;
on each mesh cell the interpolant's derivative is the constant
$40(F_k-F_{k-1})$. Errors for \(P_F\) from the chord through its starred samples are at most
\[
\begin{array}{c|ccc}
 F&\mathsf K&q&q^2\\\hline
 \text{error}&.000051&.000128&.000076
\end{array}
\]
Indeed sample errors are bounded by \(G_F/12800\) for \(G_F=\sup |\ddot F|\); also
\(|\ddot P_F|\le G_F+2L_F+o_F\) where \(L_F=\sup|\dot F|,\ o_F=\operatorname{osc} F\). Differentiating the quotient, the terms besides the weighted mean of \((s')^2\ddot F\) are
\(2\operatorname{Cov}(s'\dot F,s'\tanh(ts'))+\operatorname{Cov}(F,(s')^2)
-2\mathbb E_*[s'\tanh(ts')]\operatorname{Cov}(F,s'\tanh(ts'))\)
in its weighted expectation \(\mathbb E_*\). Bounding covariances by one quarter of the products of their ranges
establishes this derivative estimate. Adding the error at the starred
nodes and the chord error of the exact average gives the total bound
$(2G_F+2L_F+o_F)/12800$.
For $(F,G_F,L_F,o_F)=(\mathsf K,.238,.066,.0345)$ and
$(q,.53,.193,.1784)$ this is below the first two stated errors.
For $q^2$ we can use \(G_F=.347,\ L_F=.099,\ o_F=.060\).
Chord errors for \(\mathsf K,\mathsf D,\mathsf C,q\) are bounded by \(.000019,.000020,.000062,.000042\) using the one-variable estimates.

Here are the sample bounds for this part. In the formulas inside the table use sampled values at \(t_i\), including the starred samples for \(P_F\). Put \(q_M=.2559\), and define
\[
\begin{aligned}
 B_0(t)&=P(t)-(\mathsf K(t)+\mathsf K(0))/2+
       d_*^K(\mathsf D(t)-\mathsf D(0))^2+d_*^C(\mathsf C(t)-\mathsf C(0))^2,\\
 B_1(t)&=-\tfrac12(\mathsf K(t)-\mathsf K(0))
     +d_*^K(\mathsf D(t)-\mathsf D(0))_+(\mathsf D(t)+\mathsf D(0)+.0214)_+,\\
 Z_1(t)&=\tanh t\,(.0178-P(t)).
\end{aligned}
\]
\[
\begin{array}{c|l|r}
 i&\text{expression}&\text{bound}\\\hline
 0{:}144&\mathsf D-\mathsf D(0)&[-10^{-8},.0314]\\
  &P&\le .01767\\
  &B_0&\le .01773\\
  &B_0-.0146 t_i^2&\le .00001\\
  &\dot{\mathsf D}&\ge -.0213\\\hline
 0{:}56&B_1+\tfrac12(P-.0049)_+&\le .000001\\
  &Z_1 &[0,.00742]\\
  &B_1+78Z_1^2 &\le .000001\\
  &P_{q^2}-(q+q_M)P_q+(q^2+q_M^2)/2&\le .01587\\\hline
 0{:}144& -\mathsf C,\quad q,\quad \min(\mathsf D,-.01065),\quad \min(P,.0049)
      &\text{nondecreasing (samples)}
\end{array}
\]
The arithmetic for these tables proceeds in the following order.
At each node $t_i=i/40$, first use the one-variable sample rules to
enclose $\mathsf K_i$, $q_i$ and
$\mathsf C_i=-d_i+6.5g''(s\sinh t_i)$, together with their indicated
derivatives. Form $\mathsf D_i=\mathsf K_i+.021\mathsf C_i$.
For the threshold table the four columns are
$\dot{\mathsf K}_i+.021\dot{\mathsf C}_i$,
$-\dot{\mathsf C}_i$, $-n_{\mathsf K,i}$ and $n_{\mathsf C,i}$.

The starred primitive samples are not point evaluations. For each
of them, start the numerator sum in its defining formula at zero and
add the profile differences times the corresponding differences of
$\cosh t_i$. Equivalently, use $a_{t_i}=s\cosh t_i$ and
$x_i=s\sinh t_i$; the common factor $s$ cancels in the quotient.
With these starred samples and the point-value columns in hand,
substitute into $B_0,B_1,Z_1$ and the endpoint-distance expression.
Squares are the full ranges of squares of the interval arguments.
At $t=0$, take $\tanh t=0$ exactly.

For the last table row, compare successive nodes $i,i+1$ for
$0\le i<144$. The differences for $-\mathsf C$ and $q$ are
bounded below by $.000005$ and $.000012$, respectively. For the two
clipped columns the lower bound is zero; clipping is applied to the
interval endpoints before taking these differences.
These operations, with the finite enclosure rules from the
one-variable table, specify both tables here.

The monotonicity line and the chord errors prove the two claimed almost-monotonicity bounds, and for \(0\le t\le y\le3.6\),
\[
 \begin{aligned}
 P(y)&\ge\min(P(t),.0049)-2\cdot.000051,\qquad P\le.0178,\\
 \mathsf D(y)&\ge\min(\mathsf D(t),-.01065)-.000040 .
 \end{aligned}
\]
To justify this implication, let $J_F$ be the piecewise affine interpolant
of the relevant samples. If $\min(F_i,c)$ is nondecreasing, then
$\min(J_F,c)$ is nondecreasing as well: each cell connects its two
clipped endpoint values monotonically, and a cell with both endpoints
above $c$ remains above $c$. Since clipping is 1-Lipschitz,
$|F-J_F|\le\epsilon$ implies, for $t\le y$,
\[
 F(y)\ge\min(F(y),c)
 \ge\min(J_F(y),c)-\epsilon
 \ge\min(F(t),c)-2\epsilon.
\]
For starred P samples the same argument uses their total chord error.
Saturated interval values in the finite calculation are the exact
constant $c$; uncertain overlapping intervals are not used to infer
monotonicity. Also \(B_0(t)\le\min(.018,.0146t^2+.00010)\) by convexity in its profile arguments. Here the added errors are \(<.000075\) (use the sample ranges for \(\mathsf D-\mathsf D(0),\mathsf C-\mathsf C(0)\)), and \(.0146t^2\) has chord error less than .000003.

\smallskip
\noindent\textbf{The primitive budget.}\quad
Take $t=|l|$. Since $\mathsf K$ is even, its primitive in physical
length is represented by $s\sinh(t)P(t)$ on the positive half-axis.
Subtracting the primitives at the two endpoints gives
$\overline{\mathsf K}=P(y)+\theta(P(y)-P(t))$, where
$\theta=\sinh l/(\sinh y-\sinh l)$.
After subtracting $B_0(y)$ from $e_K+\Gamma$, the remaining terms
are $\theta(P(y)-P(t))$ and
\[
 -\tfrac12(\mathsf K(t)-\mathsf K(0))+
 d_*^K (\mathsf D(t)-\mathsf D(0))(\mathsf D(t)+\mathsf D(0)-2\mathsf D(y))
 +d_*^C (\mathsf C(t)-\mathsf C(0))(\mathsf C(t)+\mathsf C(0)-2\mathsf C(y)).
\]
We first bound this common correction, before distinguishing the sign
of $l$. Put $X=\mathsf D(t)-\mathsf D(0)$ and
$\varepsilon=.000040$. The product with coefficient $d_*^K$ is
\[
 X\bigl(\mathsf D(t)+\mathsf D(0)-2\mathsf D(y)\bigr)
 =X\bigl(2[\mathsf D(t)-\mathsf D(y)]-X\bigr).
\]
The node and chord bounds give $-.000021\le X\le.032$.
If $X<0$, use $\mathsf D(y)-\mathsf D(0)\le.032$ to bound
this product above by $.000021(2\cdot.032+.000021)$.
If $X\ge0$ and $\mathsf D(t)\le-.01065$, clipped monotonicity
gives $\mathsf D(y)\ge\mathsf D(t)-\varepsilon$. Hence
\[
 X\bigl(2[\mathsf D(t)-\mathsf D(y)]-X\bigr)
 \le -X^2+2\varepsilon X\le\varepsilon^2.
\]
In the other case, $\mathsf D(t)>-.01065$, the same clipped bound
gives $\mathsf D(y)\ge-.01065-.000040>-.0107$.
The product is then bounded by
$(\mathsf D(t)-\mathsf D(0))_+
 (\mathsf D(t)+\mathsf D(0)+.0214)_+$,
the product already present in $B_1(t)$.

For the $\mathsf C$ product, suppose first
$\mathsf C(t)>\mathsf C(0)$. Almost-monotonicity bounds the first
factor by $.000124$, while the global range bounds the second factor
above by $.320$. If $\mathsf C(t)\le\mathsf C(0)$, set
$a=\mathsf C(0)-\mathsf C(t)\ge0$.
The bound $\mathsf C(y)\le\mathsf C(t)+.000124$ makes the product
at most $-a^2+2(.000124)a\le.000124^2$.
After multiplying these errors by $d_*^K$ and $d_*^C$, their sum is
less than $.000014$. Thus the common correction is at most
$B_1(t)+.000014$.

We next pass from the samples of $B_1$ to its true values. The product
$(x-b)_+(x-b')_+$ is zero up to the larger threshold and is a
nondecreasing convex quadratic thereafter. Thus the product in $B_1$
is convex in $\mathsf D(t)$, with $\mathsf D(0)$ fixed.
Its two clipped factors sum to at most $.058$, both for the chords
and for the true arguments. Here we use $\mathsf D(0)<-.0140$;
the exact formula is
$(1-6.5t_c)((2+\pi)c_*-1)-2t_cc_*$, with $c_*=\log2-1/2$.
Consequently the upper error of $B_1$ from the chord through its
sample values is at most
$\tfrac12\cdot.000019+d_*^K\cdot.058\cdot.000020<.000020$.

If $l<0$, then $\theta\in[-1/2,0]$.
The lower bound
$P(y)\ge\min(P(t),.0049)-2\cdot.000051$ therefore gives
$\theta(P(y)-P(t))\le.5(P(t)-.0049)_++ .000051$.
Use the table row for $B_1+\tfrac12(P-.0049)_+$.
Convexity of the positive part and the chord error of $P$ add at
most $.000051/2$ to the $B_1$ chord error. The total correction is
bounded by
\[
 .000014+.000020+.000001+.000051/2+.000051<.00015.
\]
Combining this with $B_0(y)\le\min(.018,.0146y^2+.00010)$ and
$y\le2h$ proves the $l<0$ line of Equation~\eqref{eq:s1}.

For $l\ge0$, $t=l$ and
$0\le\theta\le\tanh l/\sinh(2h)$.
Since $P(y)\le.0178$ and $P(t)\le.0178$, the nonnegative function
$Z_1(t)=\tanh t(.0178-P(t))$ satisfies
$\theta(P(y)-P(t))\le Z_1(t)/\sinh(2h)$.
This time we use the table row for $B_1+78Z_1^2$.

The error of $Z_1$ from its starred chord is at most $.000066$.
To see this, combine the separate errors in the two factors with
$\tfrac14|\Delta P_i^*||\Delta\tanh t_i|$, the difference between
the product of their chords and the chord of their products.
The derivative estimate $|\dot P|\le L_{\mathsf K}+o_{\mathsf K}/4$
follows by the same weighted differentiation used above. Explicitly,
\[
 .000051+\frac{.0178+.007+.238/12800}{12800}
 +\frac1{4\cdot40}\left((.066+.0345/4)/40+2\cdot.238/12800\right)<.000066 .
\]
The starred samples lie in $[0,.00742]$. Convexity of the square
therefore bounds the error in $78Z_1^2$ by
$78\cdot.000066(2\cdot.00742+.000066)$.
Including the common correction, the $B_1$ chord error and the table
margin gives
$ .000014+.000020+.000001+
78\cdot.000066(2\cdot.00742+.000066)<.00015$.
It follows that
\[
 e_K+\Gamma-B_0(y)\le .00015+Z_1(t)/\sinh(2h)-78 Z_1(t)^2 ,
\]
Maximizing $Z/\sinh(2h)-78Z^2$ over real $Z$ gives
$1/(312\sinh^2(2h))$. Together with the bound on $B_0$, this proves
\eqref{eq:s1} when $l\ge0$, completing the primitive estimates.

\smallskip
\noindent\textbf{The q endpoint-distance bound.}\quad
Finally, on the part \(|l|\le t'\le y\), \(q(t')\) lies between the two endpoint values up to .000084, so the half sum of squared differences pointwise is at most
\(.1784^2/2+.000084( .1784+.000084)<.0163\). If $l<0$, reflect the extra part $[l,|l|]$ to $[0,t]$, where $t=|l|$.
To bound its mean endpoint-distance, maximize the convex quadratic in
the other endpoint value $q(y)\in[q(t)-.000084,q_M]$.
The upper endpoint $q_M$ suffices: throughout $0\le t\le1.4$,
the $q_{\max}$ column in the proof of Lemma~\ref{lem:scalar-layer}
(rows through $40{:}68$) gives $q(t)\le.245$, whereas
$q(v)\le q(t)+.000084$ on $[0,t]$. Hence
$q_M-q(t)>3(.000084)$, so the squared distance to $q_M$ is at least
the squared distance to $q(t)-.000084$ at every such $q(v)$.
After averaging, we obtain the following conservative bound:
\[
 P_{q^2}(t)-(q(t)+q_M)P_q(t)+(q(t)^2+q_M^2)/2+.000084^2/2
 \qquad (t=|l|).
\]
The table applies with upper error at most .00018 on its profile expression, proving \(b_a\le.0163\). Indeed for the product use the given chord errors plus \(\tfrac14|\Delta P_{q,i}^*||\Delta q_i|\), with \(|\dot P_q|\le L_q+o_q/4\); altogether the upper error is at most
\[
\begin{aligned}
 &.000076+2q_M\cdot.000128+(q_M+.53/12800)\cdot.000042\\
 &\quad+\frac{.193}{4\cdot40}
       \left((.193+.1784/4)/40+2\cdot.53/12800\right)\\
 &\quad+\tfrac12\cdot.000042(2q_M+.000042)<.00018,
\end{aligned}
\]
using also convexity for the \(q(t)^2/2\) term.
This completes the proof.
\end{proof}

\subsubsection{Consequences for weighted chord errors}

We first suppose $[l,y]\subset[0,3.6]$ and apply the concentration
bound to $F=(A_j)_+$. The kernel $1-\beta^2$ has the layer-cake identity
\[
 \int_{-1}^1(1-\beta^2)F(m+h\beta)\,\frac{d\beta}{2}
 =\int_0^1 2(u')^2
   \left(\frac1{2u'}\int_{-u'}^{u'}F(m+h\beta)\,d\beta\right)du'.
\]
The expression in parentheses is the uniform mean on the interval
$[m-hu',m+hu']$, whose length is $2hu'$.
For $hu'\le T_0$ it is at most $\alpha_j-\beta_jhu'$.
For longer intervals it is still at most $\alpha_j-\beta_jT_0$:
choose a subset of length $2T_0$ with at least as large a mean and
apply the concentration bound to that subset.
Since $\min(hu',T_0)\ge u'\min(h,T_0)$, integration gives
\[
 \int_0^1 2(u')^2
       [\alpha_j-\beta_j u'\min(h,T_0)]\,du'
 =\frac{2\alpha_j}{3}-\frac{\beta_j\min(h,T_0)}2
 =L_j(\min(h,T_0)).
\]

For a segment crossing zero we need only $F=(n_{\mathsf C})_+$,
which is even. Reflect the interval $[m-hu',m+hu']$ to $[0,3.6]$.
The pushforward of its length measure has density
$\mathbf1_{[m-hu',m+hu']}(t)+
\mathbf1_{[m-hu',m+hu']}(-t)$, which is at most two and has
mass $2hu'$. Divide this density by two. It is now bounded by one
and has mass $hu'=2(hu'/2)$, so the density version of the
concentration estimate bounds the original interval mean by
$\alpha_C-\beta_Chu'/2$.
Here $h/2\le T_0$ ensures that every parameter $hu'/2$ is allowed.
The same layer-cake integral therefore gives $L_C(h/2)$, with no
extra factor of two.

Apply these bounds to the chord identity and use
$W(\beta)\le g_h(1-\beta^2)$. They give
\[
 -c_a\le h^2 g_h L_C(h_s')/3,\qquad
 h_s'=\begin{cases}\min(h,T_0)& l\ge0,\\ h/2&l<0,\ h/2\le T_0.\end{cases}                 \tag{s2}\label{eq:s2}
\]
For $l\ge0$ and $.16\le h\le T_0$, the chord identity with
$F=\mathsf K$ similarly gives
$e_K\le h^2g_hL_N(h)/3$.
The endpoint differences in $\Gamma$ use the unweighted means.
Put $d_s=\alpha_D-\beta_Dh$ and $b_s=\alpha_E-\beta_Eh$.
The mean of $\dot{\mathsf D}$ is between $-.022$ and $d_s$,
and $d_s>.022$ on this width range. Therefore
$|\mathsf D(y)-\mathsf D(l)|\le2hd_s$.
For $-\dot{\mathsf C}$, almost-monotonicity supplies the lower
mean bound $-.000124/(2h)$, while concentration supplies the upper
bound $b_s$. Since $b_s>.000124/(2h)$, this gives
$|\mathsf C(y)-\mathsf C(l)|\le2hb_s$, and hence $|c_d|\le hb_s$.
Substituting these two endpoint differences in $\Gamma$ proves
\[
 e_K+\Gamma\le h^2\big[g_h L_N(h)+12d_*^K d_s^2+12d_*^C b_s^2\big]/3.                \tag{s3}\label{eq:s3}
\]
We abbreviate \(p=p_a, f=f_0^*, \chi=t_+\) in the rest of this segment proof, and put
\[
 G_0=r^2w^2,\qquad d_h=2r^2 B_h\sin(wh),\quad p_d=-d_h\tanh m.
\]
We will use the signs of the two half-differences as well as their
absolute bounds. Angular thresholds such as $\pi/w$ use the usual
circular constant. For $h\ge.16$, the bound on $|B_h\sin(wh)|$
and $\coth(.16)\le1/.16+.16/3$ give $|d_h|<.031$.
If $h\le\pi/w$, then $d_h\ge0$: use $\coth h\ge1/h$ and
$\sin z-z\cos z\ge0$ for $0\le z\le\pi$.
In this case $-.031\le p_d\le0$, while range and almost-monotonicity give
$-.08\le c_d\le.000062$. Therefore
$-.08\le c_d-p_d\le.031062$, proving $|c_d-p_d|\le.08$.
When $l\ge0$ and $.16\le h\le T_0$, use the sharper lower
bound $c_d\ge-hb_s$ instead. Since $hb_s>.032>.031062$,
we obtain $|c_d-p_d|\le hb_s$.

\subsection{Middle widths}\label{subsec:segments-middle}

We remain in the bounded region \(m\ge0\), \(|l|\le1.4\),
\(y\le3.6\), and now take \(.16\le h\le T_0\).
We first establish $p/h^2>.33$.
If $wh\le\pi/2$, use $|A_h|\le1$ and
$1-\cos(wh)\ge4(wh)^2/\pi^2$, the latter following from
monotonicity of $\sin z/z$.
For $\pi/2\le wh\le wT_0<5\pi/6$, we have
$\sin(wh)\ge1/2$, $\cos(wh)\ge-1$, and $w\coth h\ge w=4.6$;
hence the numerator of $A_h$ is positive. Thus
$p\ge2r^2$ and $p/h^2\ge2r^2/T_0^2>.33$.
The chord estimate now gives
$c_a\le.79h^2g_h/(1.5\cdot3)<p$.

First treat $l<0$. The endpoint-distance moment bound, valid through
$h=.54$, gives $b_a/h^2\le B_a^*=(.193)^2(4/3+.05h^2)$.
For the half-difference, combine $|q_+-q_-|\le2(.193)h$ with
the chord error
$|\bar q-(q_++q_-)/2|\le.53h^2g_h/(1.5\cdot3)<h^2/5$.
Thus $|b_d|\le.078h^3=:B_d^*$.
By Equations~\eqref{eq:s1} and~\eqref{eq:s2}, it suffices for Equation~\eqref{eq:s0} that
\[
 (1+\chi)f>
 .0584 h^2+.00025+\chi(p+h^2 B_a^*)+
 \frac{(p+h^2(g_h L_C(h/2)/3+B_a^*))^2+(.08+B_d^*)^2}{4\lambda}.                  \tag{s4}\label{eq:s4}
\]

\subsubsection{Nonnegative endpoints and the q variance}
For $l\ge0$ in this range put
\[
 a_s=p/h^2+g_h L_C(h)/3,\quad z=|q_+-q_-|/(2h)\le .193,\quad v_h=h\coth h-h.
\]
Then \(0\le p-c_a\le h^2 a_s\). To lower bound \(V_q\) use that the averaging density with respect to \(dt\) on \([l,y]\) is at least \(v_h/(2h)\) (since \(\cosh(t)/\cosh m\ge e^{-h}\)). No monotonicity of $q$ is needed. For a nonnegative continuous function
$\psi$, the derivative bound and the fundamental theorem of calculus give
\[
 \int_l^y\psi(q(t))\,dt
 \ge\frac1{.193}\int_l^y\psi(q(t))|\dot q(t)|\,dt
 \ge\frac1{.193}\left|\int_{q_-}^{q_+}\psi(v)\,dv\right|.
\]
Applied to the test \((q-\bar q)^2\) (with \(\bar q\) fixed and the endpoint-value interval of length \(2hz\), where the integral of this squared deviation is at least \((2hz)^3/12\)), this gives
\[
 V_q/h^2\ge v_h z^3/(3\cdot .193).
\]
We have \(0\le b_a-V_q\le h^2(z^2+h^2/25)\) (it is the square of the chord-mean error plus the squared half difference) and \(|b_d|\le2h^3 z/5\). For Equation~\eqref{eq:sv} divided by \(h^2\) we use
\[
 \begin{aligned}
 &\frac{(c_a-p-b_a)^2+(c_d-p_d-b_d)^2}{4\lambda h^2}
       +\frac{\chi(p+b_a)-(1+\chi)V_q}{h^2}\\
 &\qquad\qquad\le \frac{h^2 a_s^2+b_s^2}{4\lambda}+\chi p/h^2+R ,
 \end{aligned}
\]
where
\[
 R= Q_s(z^2+h^2/25)+h^2 b_s z/(5\lambda)+h^4 z^2/(25\lambda)
       -\frac{.855}{3\cdot .193}v_h z^3 ,
 \qquad Q_s=\chi+h^2(2a_s+.05)/(4\lambda).
\]
In fact in the first square use
\[
 (p-c_a+b_a)^2\le(h^2a_s+b_a-V_q)^2+2(h^2a_s+.0163)V_q .
\]
Here $z^2+h^2/25\le.193^2+.54^2/25=.048913<.05$.
The upper bound $h^2a_s\le.16025$ proved below gives
\[
 1-\frac{h^2a_s+.0163}{2\lambda}
 \ge 1-\frac{.16025+.0163}{1.3}>.86419>.855.
\]
More explicitly, put $D_q=b_a-V_q$, $H_a=h^2a_s$, and
$\delta_q=z^2+h^2/25$. Then $0\le D_q\le h^2\delta_q$ and
\[
\begin{aligned}
 (p-c_a+b_a)^2
 &\le(H_a+D_q+V_q)^2\\
 &=(H_a+D_q)^2+2(H_a+b_a)V_q-V_q^2\\
 &\le(H_a+D_q)^2+2(H_a+.0163)V_q.
\end{aligned}
\]
The other square is at most $(hb_s+2h^3z/5)^2$, while
$\chi b_a-(1+\chi)V_q=\chi D_q-V_q$.
After division by $h^2$, the terms involving $D_q$ are bounded by
$Q_s\delta_q$, since $\delta_q\le.05$. Expanding the other square
produces $h^2b_sz/(5\lambda)+h^4z^2/(25\lambda)$ beyond its constant
term. The retained coefficient of $-V_q/h^2$ is at least $.855$;
substitution of the variance lower bound gives precisely $R$.

\subsubsection{A cubic bound for the remaining variance cost}
We verify that \(R/h^2\le .027\), by setting \(Z=z/h\). An upper bound on each interval \([L,H]\) with successive endpoints \(.16,.23,.29,.35,.40,.45,.50,.54\) is the maximum on \(Z\ge0\) of \(Q/25+A Z^2+B Z-N_0 Z^3\). The following table bounds $h^2a_s,Q,A,B$ upwards and $N_0$ downwards.
Its final column bounds the maximum of the cubic formed from those
rounded coefficients. Every entry is scaled by $10^5$; the seven rows
correspond to the seven successive width intervals. The columns are
rounded independently from their defining formulas with the unrounded
$A_0$; a column is not calculated recursively from an earlier rounded
column.
\[
\begin{array}{rrrrrr}
 h^2a_s&Q&A&B&N_0&\sup(R/h^2)\\\hline
8298&12885&12902&1576&20047&2028\\
13121&16655&16699&1851&26748&2449\\
15242&18360&18453&2069&31598&2526\\
15218&18414&18572&2208&35687&2308\\
15449&18673&18925&2309&38555&2242\\
15768&19010&19395&2370&40970&2218\\
16025&19288&19811&2390&42963&2195
\end{array}
\]
For details put
\[
 S_u(h)=1+h^2/6+.0084h^4,\quad H_u(h)=1+h^2/2+.0421h^4,\qquad
 g_u(h)=.75 S_u(h)(1+H_u(h)).
\]
These are upper bounds for $S$, $\cosh h$, and $g_h$ on the present range.
For example, after dividing the remainder of $S$ by $h^4$, its first
term is $1/120$ and its successive term ratio is at most $.54^2/42$.
Thus its sum is at most $875/104271<.0084$. For $\cosh h$ the
corresponding bound is $3125/74271<.0421$. The product defining $g_u$
has nonnegative factors, so it also has the claimed direction. For the first column use
\[
 A_0=g_u(H)H^2L_C(H)/3+
 \min\left(G_0 H^2(4/3+.05H^2),\ 2r^2\left[1-\frac{1-\sqrt{1+w^2(1/L+L/3)^2}}{2(1+w^2)}\right]\right).
\]
Indeed \(h^2L_C(h)\) is increasing here and we bounded \(p\) as in the short-interval and tail estimates, using \(\coth h\le1/h+h/3\). Use
\[
\begin{aligned}
 Q&=\chi+\frac{2A_0+.05H^2}{4\lambda},&
 A&=Q+\frac{H^4}{25\lambda},\\
 B&=\frac{H(\alpha_E-\beta_EH)}{5\lambda},&
 N_0&=\frac{.855L(1-L+L^2/3-.023L^4)}{3\cdot.193}.
\end{aligned}
\]
The endpoint choices have the required directions.
The derivatives of $h^2L_C(h)$ and $hb_s$ are positive through $.54$;
for the latter, $(hb_s)'=.281-.508h\ge.00668$.
Also $hv_h=2h^2/(e^{2h}-1)$ is increasing: its derivative has the sign
of $(1-h)e^{2h}-1$. The function $(1-h)e^{2h}$ increases to $h=.5$
and then decreases, while its value at $.54$ is greater than
$.46(1+1.08+1.08^2/2)>1$.
Finally, the coefficient of $h^{2n}$, $n\ge2$, in
$\cosh h-(1+h^2/3-.023h^4)S$, multiplied by $(2n-3)!$, is
$.023-1/[3(4n^2-1)]>0$. The lower coefficients vanish, proving
$h\coth h\ge1+h^2/3-.023h^4$ and hence the lower bound $N_0$. The maximum for the rounded bounds in the display uses \(Z_*=(A+\sqrt{A^2+3N_0 B})/(3N_0)\) and is \(Q/25+(A Z_*^2+2B Z_*)/3\). The formulas generating these bounds use rational operations and two
positive square roots, in $A_0$ and $Z_*$. For a positive rational $a$, a fully rational
rule is to put $k=\lfloor\sqrt{\lfloor10^{160}a\rfloor}\rfloor$ by
integer square comparison and use
$\sqrt a\in[k/10^{80},(k+1)/10^{80}]$. Substitution with outward
rational arithmetic verifies the stated coefficient bounds. Using the rounded coefficients gives, in row order, cubic
maxima strictly below
\[
 .02028,\ .02449,\ .02526,\ .02308,\ .02242,\ .02218,\ .02195.
\]
In the fourth row the displayed upper bound $B=.02208$ is exact;
all rounding inequalities here are non-strict. Every cost bound is below
$.027$, which proves the required uniform estimate.

Consequently, Equation~\eqref{eq:s3} reduces Equation~\eqref{eq:sv} to
\[
 (1+\chi) f/h^2 >
 \big[g_h L_N(h)+12d_*^K d_s^2+12d_*^C b_s^2\big]/3+
 (h^2 a_s^2+b_s^2)/(4\lambda)+\chi p/h^2+.027h^2.                                 \tag{s5}\label{eq:s5}
\]

\subsubsection{Exact polynomial certificates}
We now prove Equations~\eqref{eq:s4} and~\eqref{eq:s5} throughout their common
width interval by finite polynomial certificates. Multiply both inequalities (right terms brought to the left) by \(S^2\). With \(\operatorname{sinc} z=(\sin z)/z\) the angular functions obey
\[
\begin{aligned}
 pS/h^2 &= \frac{2G_0}{1+w^2}
  \big[(S-\cosh h\,\operatorname{sinc}(2wh))/h^2+S\,\operatorname{sinc}^2(wh)\big],\\
 f S^2/h^2 &= \frac{G_0}{1+w^2}(S^2-\operatorname{sinc}^2(wh))/h^2 .
\end{aligned}
\]
The coefficient sequences below represent the four even series
\[
 S=\sum_{i\ge0}s_i h^{2i},\qquad
 \cosh h=\sum_{i\ge0}a_i h^{2i},\qquad
 \operatorname{sinc}(2wh)=\sum_{i\ge0}d_i h^{2i},\qquad
 \operatorname{sinc}^2(wh)=\sum_{i\ge0}e_i h^{2i},
\]
where $s_i=1/(2i+1)!$, $a_i=1/(2i)!$,
$d_i=(-4w^2)^i/(2i+1)!$, and $e_i=2(-4w^2)^i/(2i+2)!$.
Ordinary convolution on nonnegative indices is denoted by $*$.
Thus $(a*d)_i$ and $(s*e)_i$ are the coefficients of the two
products in the formula for $pS/h^2$.
We use the following upper and lower polynomials:
\[
\begin{aligned}
 pS/h^2\ \le\ P_u&=.00025+\frac{2G_0}{1+w^2}
   \left[\sum_{i=1}^7(s_i-(a*d)_i)h^{2i-2}+\sum_{i=0}^7(s*e)_i h^{2i}\right],\\
 f S^2/h^2\ \ge\ F_*&=\frac{G_0}{1+w^2}\sum_{i=1}^6
         \frac{2\cdot4^i(1-(-w^2)^i)}{(2i+2)!}h^{2i-2}.
\end{aligned}
\]
For $F_*$, the omitted series starts at $i=7$. It alternates with
first term positive and decreasing magnitudes: the successive ratios
from that index are bounded by $8w^2h^2/(18\cdot17)<1$.
Its sum is therefore nonnegative, which gives the lower bound.
For $P_u$, both omitted series start at $i=8$.
Before the common multiplier, their absolute coefficients are bounded by
\[
 \frac{1+9.26^{2i+1}/(2w)}{(2i+1)!},\qquad
 \frac{1+(1+2w)^{2i+3}}{2w^2(2i+3)!}.
\]
These bounds follow by expanding $\cosh h\sin(2wh)$ and
$\sinh h\cos(2wh)$ in complex exponentials, whose frequencies
have modulus $\sqrt{1+4w^2}<9.26$.
Multiply the first coefficient bound by $h^{2i-2}$, the second by
$h^{2i}$, and both by $2G_0/(1+w^2)$.
At $i=8$, setting $h=.54$ makes their sum less than $.000153$.
For every later term the ratio is below $.10$: the factorial
denominator factors increase with $i$, whereas the geometric
numerator ratios and $h^2\le.54^2$ stay bounded.
The entire omitted tail is consequently less than
$.000153/(1-.10)<.00025$, the allowance included in $P_u$.

The following polynomials lower bound the two deficits after multiplication
by $S^2$, in the order of Equations~\eqref{eq:s4} and~\eqref{eq:s5}.
Indeed, $P_u\ge pS/h^2\ge0$, and all other factors occurring inside
the squares are nonnegative. The linear factors also have the required
signs: on this interval $L_N(h)\ge.08521>0$ and
$L_C(h)>.26609>0$. Replacing $S$ by $S_u$ and $g_h$ by $g_u$
therefore increases every upper cost and gives
\[
\begin{aligned}
 T_{(4)}={}&(1+\chi)h^2F_*- S_u^2\big[.0584h^2+.00025+\chi h^2 B_a^*+(.08+B_d^*)^2/(4\lambda)\big]\\
          &-\chi h^2 P_u S_u-h^4(P_u+S_u(g_u L_C(h/2)/3+B_a^*))^2/(4\lambda),\\
 T_{(5)}={}&(1+\chi)F_*-
  S_u^2\left[\big(g_u L_N(h)+12d_*^K d_s^2+12d_*^C b_s^2\big)/3+b_s^2/(4\lambda)+.027h^2\right]\\
          &-h^2(P_u+S_u g_u L_C(h)/3)^2/(4\lambda)-\chi P_u S_u .
\end{aligned}
\]
Both polynomials are positive on $[.16,.54]$, as the following finite
rational certificate shows. Put $h=.16+.38u$, $0\le u\le1$, and write
\[
 T(.16+.38u)=\sum_{r=0}^{d}\beta_r\binom dr u^r(1-u)^{d-r},
 \qquad d=\deg T.
\]
The basis functions are nonnegative and sum to one. It therefore
suffices that every $\beta_r$ be strictly positive. The degrees are 32,30 respectively. Lower bounds for coefficients multiplied by \(10^5\), grouped by index \(0{:}7,8{:}15,16{:}23,24{:}32\) (through 30 in row two) are
\[
\begin{array}{c|rrrr}
 &0{:}7&8{:}15&16{:}23&24{:}d\\\hline
 T_{(4)}&173&576&770&412\\
T_{(5)}&1586&701&408&407
\end{array}
\]
To verify by sums, if \(t_j=[h^j]T\) from the displayed polynomials, take for each index \(r\) the sum
\[
 \sum_{i=0}^r\frac{\binom ri}{\binom{\deg T}i}(.38)^i
       \sum_{j=i}^{\deg T}t_j\binom ji(.16)^{j-i}.
\]
The expression is a finite rational sum: the coefficients $t_j$ are
specified by the preceding products, factorials, and convolutions.
Binomial expansion first gives the power coefficients in $u$; the identity
$u^i=\sum_{r=i}^d\binom ri\binom di^{-1}\binom dr u^r(1-u)^{d-r}$
then gives the displayed basis change. Direct substitution yields the
grouped lower bounds. Since each is positive, the Bernstein
representation proves strict positivity on the entire closed interval. This completes the middle widths through .54.

\subsection{Larger widths}\label{subsec:segments-large}

We remain in the bounded region \(m\ge0\), \(|l|\le1.4\),
\(y\le3.6\), and now take \(h\ge.54\).
Equation~\eqref{eq:s1} controls \(e_K+\Gamma\) throughout this region;
write \(b_M=.0163\). For $.54\le h\le\pi/w$, put
$E(h)=\cos^2(wh)+\tfrac12w\coth h\sin(2wh)$.
Since $wh\in(3\pi/4,\pi]$, we have $\sin(2wh)\le0$ and
$\cos(2wh)\ge0$. Thus
\[
 E'(h)=-w\sin(2wh)-\tfrac12w\operatorname{csch}^2h\sin(2wh)
          +w^2\coth h\cos(2wh)\ge0.
\]
As $p=2r^2(1-E/(1+w^2))$, this proves that $p$ is decreasing. Thus \(p>.08\). Also \(f\) is increasing up to \(\pi/w\). In any part of the current range \(-c_a\le .0802\), by almost-monotonicity (all values on the segment are at least \(\mathsf C(y)-.000124\)), and \(c_a\le .160\) by the profile range. On \([L,H]\) inside or covering part of \([.54,\pi/w]\), Equation~\eqref{eq:s2} gives additional bounds on \(-c_a\). For the intervals with endpoints \(.54,.58,.63,.684\), write
\[
 R_P=\min(.0802, H^2 g_H L_C(.54)/3),\quad R_O=\min(.0802,H^2 g_H L_C(H/2)/3)
\]
(nonnegative or opposite signs of endpoints, respectively), using monotonicity of \(h^2L_C(h/2)\) here. Then \(|c_a-p-b_a|\le R_j+p+b_M\) in case \(j\in\{P,O\}\), and \(|c_d-p_d-b_d|\le .08+b_M\) since \(|b_d|\le b_a\).
Here are bounds scaled by \(10^5\), rounding \(p\) up, \(f\) down. Column \(U_P\) bounds \(e_K+\Gamma\) for \(l\ge0\); for \(l<0\) use .01815 before scaling. Last two columns bound the remaining budgets below in the two cases:
\[
\begin{array}{rrrrrrr}
 p& f&U_P&R_P&R_O& P&O\\\hline
11379&4367&2004&5141&7101&2190&1901\\
10895&4500&1970&6216&8020&2278&2006\\
10112&4593&1937&7539&8020&2348&2276
\end{array}
\]
Indeed use \(p(L),f(L)\), \(.01815+1/(312\sinh^2(2L))\), the displayed formulas and, for the budgets,
\[
 (1+\chi)f-\chi(p+b_M)-\big[(R_j+p+b_M)^2+(.08+b_M)^2\big]/(4\lambda).
\]
For these numerical entries, Taylor polynomials through degree 20
for the trigonometric and hyperbolic functions have absolute error below
$10^{-8}$: their arguments are at most three (use the angle $wL$), and
\[
 \frac{e^3 3^{21}}{21!}<\frac{21\cdot3^{21}}{21!}<10^{-8}.
\]
The bound $e^3<21$ follows, for example, from $e<11/4$.
The last grid endpoint $.684$ exceeds $\pi/w$, so the last row is used
only up to $\pi/w$; the larger endpoint provides valid radius upper
bounds throughout that row. This proves Equation~\eqref{eq:s0} on this part.

It remains to consider $h\ge\pi/w>.68$.
The angular range bound gives $p>.077$, so the existing bounds
$-.0802\le c_a\le.160$ and $0\le b_a\le b_M$ imply
$|c_a-p-b_a|\le.0802+p+b_M$.
For the other squared term, the sign of $d_h$ determines which
estimate is stronger. If $d_h\ge0$, then
$-.031\le p_d\le0$, and the bounds on $c_d$ still give
$|c_d-p_d|\le.08$.

If $d_h<0$, expand its definition as
\[
 d_h=\frac{2r^2}{1+w^2}
       \bigl(\coth h\sin^2(wh)-w\sin(wh)\cos(wh)\bigr).
\]
Replacing $\coth h$ by $1$ decreases the quadratic form in
parentheses. Its smallest eigenvalue is
$(1-\sqrt{1+w^2})/2$, so in this sign case
$|d_h|\le r^2(\sqrt{1+w^2}-1)/(1+w^2)<.00810$.
Thus $|c_d-p_d|\le.08+.00810=.08810$.
There is a compensating improvement in $p$: negativity of $d_h$
forces $\sin(wh)\cos(wh)>0$, hence
$A_h\cos(wh)>0$ and $p\le2r^2=.0968$.

For $\pi/w\le h\le.9$, the same bound $p\le.0968$ holds
without a sign condition on $d_h$. Indeed
$\pi\le wh\le4.14<3\pi/2$, so both sine and cosine are
nonpositive and $A_h\cos(wh)\ge0$.
To bound $f$ on this interval, first let $h\le.78$.
Then $|\sin(wh)|\le w(h-\pi/w)\le w(.78-.68)$, while
$\sinh h\ge\sinh(.68)$.
For $.78\le h\le.9$, use $|\sin(wh)|\le1$ and
$\sinh h\ge\sinh(.78)>.85$ instead.
Inserting either pair in the definition of $f$ gives $f\ge.0431$.
The remaining budget, with $.0802$ in place of $R_j$ and $.08810$
in place of $.08$, is therefore at least $.02005$.
The primitive cost in Equation~\eqref{eq:s1} is at most $.01912$.

For $h\ge.9$, the same variance formula with
$\sinh h\ge\sinh(.9)$ gives $f\ge.04413$.
The general angular range estimate gives $p\le.10882$, and
Equation~\eqref{eq:s1} costs at most $.01853$.
Use the improved $p$ bound when $d_h<0$, and the improved
half-difference bound when $d_h\ge0$. The two substitutions into
the remaining-budget formula are
\[
\begin{array}{c|cc|cc}
 &p\ \text{at most}&|c_d-p_d|\ \text{at most}
 &\text{budget at least}&\text{cost at most}\\\hline
 d_h<0&.0968&.08810&.02115&.01853\\
 d_h\ge0&.10882&.08&.01916&.01853
\end{array}
\]
All these numerical estimates follow, for example, from
$\coth(.68)\le1/.68+.68/3$, $\coth(.9)<1.3965$,
$\sinh(.78)>.85$, $\sinh(.9)>1.024$,
$\sinh(1.36)>1.819$ and $\sinh(1.8)>2.94$.
The positive series for $\exp(2h)$ and $\sinh h$ through degree
seven give these elementary bounds.
The three remaining budgets exceed their respective costs:
\[
 .02005>.01912,\qquad .02115>.01853,\qquad .01916>.01853.
\]
They therefore prove Equation~\eqref{eq:s0} in all the final width cases.

\paragraph{Completion of the proof of Proposition~\ref{prop:segment-estimate}.}
After reflection, every distinct endpoint pair has $m\ge0$ and $h>0$.
Subsection~\ref{subsec:segments-short} covers $h\le.16$. For larger $h$,
Subsection~\ref{subsec:segments-tails} covers $l\le-1.4$, $l\ge1.4$, and
$|l|<1.4$ with $y\ge3.6$. The remaining bounded region has
$l\in[-1.4,1.4]$ and $y\le3.6$; its widths are covered by
Subsections~\ref{subsec:segments-middle} and~\ref{subsec:segments-large}.
The boundaries $l=\pm1.4$, $y=3.6$, and
$h=.16,.54,\pi/w,.9$ each belong to at least one of these arguments.
The positive polynomial coefficients and the strict budget margins
prove strictness for every $h>0$. By the reduction in
Subsection~\ref{subsec:segments-reduction}, this is the required inequality.
At coincident endpoints every deviation and endpoint difference
vanishes, giving the asserted equality.\qed
